\documentclass[11pt]{article}

\usepackage[T1]{fontenc}
\usepackage[utf8]{inputenc}
\usepackage{lmodern}
\usepackage[margin=1in]{geometry}
\usepackage{amsmath,amssymb,amsthm,mathtools}
\usepackage{microtype}
\usepackage{graphicx}
\usepackage{booktabs,tabularx,array}
\usepackage{enumitem}
\usepackage{xcolor}
\usepackage[hidelinks]{hyperref}

\newtheorem{theorem}{Theorem}[section]
\newtheorem{proposition}[theorem]{Proposition}
\newtheorem{lemma}[theorem]{Lemma}
\newtheorem{corollary}[theorem]{Corollary}
\theoremstyle{definition}
\newtheorem{definition}[theorem]{Definition}
\theoremstyle{remark}
\newtheorem{remark}[theorem]{Remark}

\newcommand{\R}{\mathbb{R}}
\newcommand{\N}{\mathbb{N}}
\newcommand{\cof}{\operatorname{cof}}
\newcommand{\norm}[1]{\left\lVert #1\right\rVert}
\newcommand{\absnorm}[1]{\left\lvert #1\right\rvert}
\newcommand{\rhobar}{\overline\rho}
\newcommand{\deltabar}{\overline\delta}
\newcommand{\deltastar}{\overline\delta^{\,*}}
\newcommand{\wstar}{w^*}
\newcommand{\eps}{\varepsilon}
\newcommand{\asympC}{\mathrel{\lesssim}}
\newcommand{\asympG}{\mathrel{\gtrsim}}

\title{Simultaneous AUS/AUC Renormings\\
\large Audit of the proposed averaging argument, a corrected quantitative theorem, and the exact separable endpoint}
\author{Research note}
\date{June 20, 2026}

\begin{document}
\maketitle

\begin{abstract}
We examine the proposed improvement of Lemma~4.3 in
Dilworth--Kutzarova--Lancien--Randrianarivony.  The proposed tail summation is a
useful observation, but the smoothness estimate used in the packet has an
unsupported dependence on the mixing parameter when the AUS exponent
exceeds~2.
For the predual of
\(N^*+\eps M^*\), the generally valid power-type estimate supplied by duality is
\(\rhobar(t)\lesssim \eps^{-(p-1)}t^p\).  The stronger
\(\eps^{-1}t^p\) estimate used in the packet does not follow from the
stated hypotheses when \(p>2\).  Consequently, the claimed theorem for every
\(r>2q\) is justified by that proof when \(p\le 2\), but not for arbitrary
\(p\).

We prove a corrected theorem valid for every reflexive Banach space: separate
\(p\)-AUS and \(q\)-AUC renormings imply a common norm which is
\(p\)-AUS and \(r\)-AUC for every \(r>pq\).  We also prove that the whole class
of positive linear Asplund averages built from
\(N^*+\eps_nM^*\) has an intrinsic \(pq\)-barrier.

Finally, the exact endpoint has a positive answer for separable reflexive
spaces.  This follows from Proposition~3.3 and Theorem~4.1(e) of
Odell--Schlumprecht: the two separate norms give exact
\((q,p)\)-tree estimates (the lower and upper branches are allowed to be
different), and hence an embedding into a reflexive space with an FDD
satisfying exact \(1\)-\((q,p)\) estimates.  Restricting the FDD norm gives a
single \(p\)-AUS, \(q\)-AUC norm on the original space.  I found no published
resolution or counterexample for the arbitrary nonseparable case posed in the
source paper.
\end{abstract}

\section{Question and conclusions}

For a normed Banach space \(X\), write \(\rhobar_X\) and \(\deltabar_X\) for
the moduli of asymptotic uniform smoothness and convexity.  The question in
Remark~4.4 of \cite{DKLR} is:

\begin{quote}
If a reflexive Banach space admits an equivalent norm with
\(\rhobar(t)\le C t^p\) and, possibly under another equivalent norm, an
estimate \(\deltabar(t)\ge c t^q\), must it admit one equivalent norm satisfying
both estimates with the same exponents \(p\) and \(q\)?
\end{quote}

\begin{center}
\includegraphics[width=\textwidth]{figures/open_problem_crop.png}
\end{center}

The outcome of the audit is summarized below.  Here ``general'' includes
nonseparable reflexive spaces.

\begin{center}
\renewcommand{\arraystretch}{1.25}
\begin{tabularx}{\textwidth}{@{}>{\raggedright\arraybackslash}p{0.31\textwidth}
>{\raggedright\arraybackslash}p{0.20\textwidth}X@{}}
\toprule
Statement & Verdict & Reason \\
\midrule
The submitted claim: common \(p\)-AUS and \(r\)-AUC for every \(r>2q\)
& Correct for \(p\le2\); gap in general
& The mixed norm has smoothness constant of order \(n^{p-1}\), not \(n\). \\
A corrected general averaging result
& Proved here
& A common \(p\)-AUS and \(r\)-AUC norm exists for every \(r>pq\). \\
Exact preservation for separable reflexive spaces
& Positive; implicit in the literature
& Odell--Schlumprecht tree estimates and their exact FDD embedding theorem. \\
Exact preservation for arbitrary nonseparable reflexive spaces
& No resolution found
& The separable FDD theorem does not directly cover nonseparable spaces, and the averaging method has a quantitative barrier. \\
\bottomrule
\end{tabularx}
\end{center}

The correction is not cosmetic.  If \(p>2\), the estimate asserted as equation
(1) in the submitted packet does not follow from AUS/AUC duality with its
stated dependence on \(n\).  The same scaling issue occurs in the displayed
argument of Lemma~4.3 in \cite{DKLR}.  This identifies a gap in that particular
proof for large \(p\); it does not by itself disprove the theorem stated there.

\section{Definitions and elementary facts}

For \(t>0\), the asymptotic moduli of a norm \(\norm{\cdot}\) are
\begin{align*}
\rhobar_X(t)
&=\sup_{x\in S_X}\inf_{Y\in\cof(X)}\sup_{y\in S_Y}
   \bigl(\norm{x+ty}-1\bigr),\\
\deltabar_X(t)
&=\inf_{x\in S_X}\sup_{Y\in\cof(X)}\inf_{y\in S_Y}
   \bigl(\norm{x+ty}-1\bigr),
\end{align*}
where \(\cof(X)\) is the family of closed finite-codimensional subspaces of
\(X\).  We say that the norm is \(p\)-AUS if
\(\rhobar_X(t)\le C t^p\) for \(0<t\le1\), and \(q\)-AUC if
\(\deltabar_X(t)\ge c t^q\) on that interval.

We use two standard facts.

\begin{enumerate}[label=(\roman*)]
\item If \(X\) is reflexive, a norm on \(X\) is \(p\)-AUS if and only if its
dual norm is weak-star AUC of power type \(p'=p/(p-1)\), up to universal
changes of constants.  More quantitatively, the two moduli are equivalent
through Young duality; see \cite{JLPS,DKLR}.
\item AUS and AUC estimates pass to closed subspaces.  Equivalent changes of
norm alter constants, but not a fixed power exponent.
\end{enumerate}

Throughout the quantitative argument we choose equivalent norms \(N\) and
\(M\) satisfying
\begin{equation}\label{eq:NM}
 N(x)\le M(x)\le C_0N(x),\qquad
 \deltabar_N(t)\ge a t^q,\qquad
 \rhobar_M(t)\le b t^p.
\end{equation}
Then
\begin{equation}\label{eq:dual-equivalence}
 C_0^{-1}N^*(x^*)\le M^*(x^*)\le N^*(x^*).
\end{equation}

\section{Audit of the submitted argument}

The packet defines, for \(n\ge1\),
\[
 L_n^*(x^*)=N^*(x^*)+n^{-1}M^*(x^*)
\]
and lets \(L_n\) be the predual norm.  Its key smoothness assertion is
\begin{equation}\label{eq:claimed}
 \rhobar_{L_n}(t)\le K n t^p.
\end{equation}
The tail summation in the convexity half is sound once a suitable smoothness
estimate is available.  The problem is the exponent of \(n\) in
\eqref{eq:claimed}.

\subsection{The correct scaling}

Let \(p'=p/(p-1)\).  Since \(M\) is \(p\)-AUS, \(M^*\) is weak-star AUC of
power type \(p'\).  For \(0<\eps\le1\), put
\begin{equation}\label{eq:mixed-dual}
 L_\eps^*=N^*+\eps M^*,
\end{equation}
and denote its predual norm by \(L_\eps\).

\begin{lemma}[Mixed-dual scaling]\label{lem:mixed}
There are constants \(c,C>0\), depending only on the data in
\eqref{eq:NM}, such that for all \(0<\eps\le1\),
\begin{align}
 \deltastar_{L_\eps^*}(s)&\ge c\eps s^{p'}
 &&(0<s\le1),\label{eq:dualconvex}\\
 \rhobar_{L_\eps}(t)&\le C\eps^{-(p-1)}t^p
 &&(0<t\le1).\label{eq:predualsmooth}
\end{align}
Moreover,
\begin{equation}\label{eq:predual-close}
 (1+\eps)^{-1}N(x)\le L_\eps(x)\le N(x).
\end{equation}
\end{lemma}

\begin{proof}
The estimate \eqref{eq:predual-close} follows immediately from
\(N^*\le L_\eps^*\le(1+\eps)N^*\), which in turn follows from
\eqref{eq:dual-equivalence}.

We sketch the quantitative point in \eqref{eq:dualconvex}.  Let
\(L_\eps^*(x^*)=1\).  By \eqref{eq:dual-equivalence}, both \(N^*(x^*)\) and
\(M^*(x^*)\) are uniformly controlled, and \(M^*(x^*)\) is uniformly bounded
away from zero.  Use weak-star AUC of \(M^*\) to choose a weak-star closed
finite-codimensional subspace on which the \(M^*\)-part gains at least a
constant multiple of \(s^{p'}\).  Since \(X\) is reflexive, choose
\(u\in X\) which norms \(x^*\) for \(N^*\), and intersect the preceding
subspace with \(\ker u\).  On this intersection the \(N^*\)-part cannot
decrease, while the \(M^*\)-part gains \(c_0s^{p'}\).  Multiplication by the
coefficient \(\eps\) in \eqref{eq:mixed-dual} gives
\eqref{eq:dualconvex}.

The Young conjugate of \(s\mapsto c\eps s^{p'}\) is, up to a constant,
\(t\mapsto \eps^{-(p-1)}t^p\).  Quantitative AUS/weak-star-AUC duality
therefore gives \eqref{eq:predualsmooth}.
\end{proof}

Putting \(\eps=n^{-1}\) in Lemma~\ref{lem:mixed} gives
\begin{equation}\label{eq:correct-n}
 \rhobar_{L_n}(t)\le C n^{p-1}t^p,
\end{equation}
whereas the argument in the packet requires \(Cnt^p\).  When \(p\le2\), the weaker bound
\(n^{p-1}\le n\) recovers the estimate used in the packet.  When \(p>2\),
the asserted dependence on \(n\) is not supplied by the cited duality
argument.  Thus the proof has a genuine gap; this observation alone does not
show that the stronger estimate is false for every possible pair of norms.

\subsection{What survives from the packet}

If the averaging weights are \(n^{-\theta}\), the smoothness sum is
\[
 \sum_{n=1}^\infty n^{-\theta}n^{p-1}
 =\sum_{n=1}^\infty n^{p-1-\theta},
\]
so the correct condition is
\begin{equation}\label{eq:theta-smooth}
 \theta>p,
\end{equation}
not merely \(\theta>2\).  The convexity tail computed in the packet remains
\[
 t^q\sum_{n\gtrsim t^{-q}}n^{-\theta}\asymp t^{q\theta}.
\]
Thus the same architecture works whenever one can choose
\(p<\theta<r/q\), namely whenever \(r>pq\).  In particular, the submitted
\(r>2q\) statement is proved by that argument for \(p\le2\), and more
generally for an individual \(r\) satisfying \(r>pq\).

\section{A corrected general theorem}

\begin{theorem}\label{thm:corrected}
Let \(X\) be a reflexive Banach space.  Suppose that \(X\) admits an
equivalent \(p\)-AUS norm and an equivalent \(q\)-AUC norm, where
\(1<p,q<\infty\).  Then, for every \(r>pq\), \(X\) admits an equivalent norm
which is simultaneously \(p\)-AUS and \(r\)-AUC.
\end{theorem}

\begin{proof}
Fix \(r>pq\), and choose
\begin{equation}\label{eq:theta-choice}
 p<\theta<r/q.
\end{equation}
For \(n\ge1\), let \(L_n\) be the predual of
\(N^*+n^{-1}M^*\), and define
\begin{equation}\label{eq:averaged-norm}
 \absnorm{x}=\sum_{n=1}^\infty n^{-\theta}L_n(x).
\end{equation}
By \eqref{eq:predual-close}, this is an equivalent norm, and there are
constants \(A,B>0\) such that, uniformly in \(n\),
\begin{equation}\label{eq:uniform-equivalence}
 A\absnorm{x}\le L_n(x)\le B\absnorm{x},\qquad
 A\absnorm{x}\le N(x)\le B\absnorm{x}.
\end{equation}

\medskip
\noindent\emph{Smoothness.}
Let \(\absnorm{x}=1\) and \(0<t\le1\).  Choose \(m\) so large that
\begin{equation}\label{eq:tail-choice}
 Bt\sum_{n>m}n^{-\theta}\le t^p.
\end{equation}
For each \(1\le n\le m\), Lemma~\ref{lem:mixed}, homogeneity, and the
uniform equivalence \eqref{eq:uniform-equivalence} give a
finite-codimensional subspace \(E_n\) such that
\begin{equation}\label{eq:individual-aus}
 L_n(x+y)\le L_n(x)+C_1n^{p-1}t^p
\end{equation}
whenever \(y\in E_n\) and \(\absnorm{y}=t\).  The constant \(C_1\) is
independent of \(n,m,x,t\).  Put \(E=\bigcap_{n\le m}E_n\).  For
\(y\in E\) with \(\absnorm{y}=t\), use \eqref{eq:individual-aus} on the
head and the triangle inequality on the tail:
\begin{align*}
 \absnorm{x+y}
 &\le \sum_{n\le m}n^{-\theta}
       \bigl(L_n(x)+C_1n^{p-1}t^p\bigr)
      +\sum_{n>m}n^{-\theta}\bigl(L_n(x)+L_n(y)\bigr)\\
 &\le 1+C_1t^p\sum_{n=1}^\infty n^{p-1-\theta}
      +Bt\sum_{n>m}n^{-\theta}\\
 &\le 1+C_2t^p.
\end{align*}
The series converges because \(\theta>p\).  Hence \(\absnorm{\cdot}\) is
\(p\)-AUS.

\medskip
\noindent\emph{Convexity.}
Let \(\absnorm{x}=1\) and \(0<t\le1\).  By the \(q\)-AUC estimate for
\(N\), uniform equivalence, and the standard rescaling of the AUC modulus,
there are \(\eta>0\) and \(E_0\in\cof(X)\) such that
\begin{equation}\label{eq:N-gain}
 N(x+y)\ge N(x)(1+\eta t^q)
\end{equation}
for every \(y\in E_0\) with \(\absnorm{y}=t\).  Choose \(\alpha>0\) large
and an integer \(n_0\) with
\begin{equation}\label{eq:n0}
 n_0-1\le \alpha t^{-q}\le n_0.
\end{equation}
For \(n\ge n_0\), equations \eqref{eq:predual-close},
\eqref{eq:N-gain}, and \(L_n(x)\le N(x)\) give
\begin{align*}
 L_n(x+y)-L_n(x)
 &\ge \frac{N(x+y)}{1+n^{-1}}-N(x)\\
 &\ge N(x)\left(\frac{1+\eta t^q}{1+n^{-1}}-1\right)
 \ge d t^q,
\end{align*}
provided \(\alpha\) is sufficiently large.  Here \(d>0\) is independent of
\(x,t,n\), using \(N(x)\ge A\).

For each of the finitely many indices \(n<n_0\), choose a norming functional
\(f_n\) for \(x\) in the norm \(L_n\), and put
\[
 E=E_0\cap\bigcap_{n<n_0}\ker f_n.
\]
Then \(L_n(x+y)\ge L_n(x)\) for \(n<n_0\) and \(y\in E\).  Therefore,
for \(y\in E\) with \(\absnorm{y}=t\),
\begin{align*}
 \absnorm{x+y}
 &\ge \absnorm{x}+dt^q\sum_{n\ge n_0}n^{-\theta}\\
 &\ge 1+d_1t^q n_0^{1-\theta}
 \ge 1+d_2t^{q\theta}.
\end{align*}
Thus the norm is \(q\theta\)-AUC.  Since \(q\theta<r\) and
\(t^{q\theta}\ge t^r\) for \(0<t\le1\), it is also \(r\)-AUC.
\end{proof}

\begin{remark}
For \(p\le2\), Theorem~\ref{thm:corrected} is weaker than the packet's
valid conclusion \(r>2q\) only when \(p<2\): in that range the packet's use
of the coarser bound \(n\) imposes \(\theta>2\), whereas the sharp scaling
allows the better threshold \(r>pq\).  At \(p=2\), the two thresholds agree.
\end{remark}

\section{A barrier for positive linear Asplund averages}

The preceding proof can be formulated with arbitrary parameters.  Let
\(\eps_n\downarrow0\), let \(w_n>0\), and let \(L_{\eps_n}\) be the predual
of \(N^*+\eps_nM^*\).  Consider
\begin{equation}\label{eq:general-average}
 \absnorm{x}=\sum_{n=1}^\infty w_nL_{\eps_n}(x).
\end{equation}
The termwise smoothness estimate from Lemma~\ref{lem:mixed} suggests the
necessary summability condition
\begin{equation}\label{eq:summability-barrier}
 \sum_{n=1}^\infty w_n\eps_n^{-(p-1)}<\infty.
\end{equation}
This condition forces a quantitative loss in the convexity tail.

\begin{proposition}[Averaging barrier]\label{prop:barrier}
Assume \eqref{eq:summability-barrier}, and define
\[
 W(s)=\sum_{\eps_n\le s}w_n.
\]
Then
\begin{equation}\label{eq:little-o}
 W(s)=o(s^{p-1})\qquad(s\downarrow0).
\end{equation}
Consequently, the AUC gain furnished by the tail mechanism at scale
\(s=t^q\) is
\begin{equation}\label{eq:barrier-gain}
 t^qW(t^q)=o(t^{pq}).
\end{equation}
In particular, no argument based solely on a positive linear average of the
mixed norms and on summing their individual tail gains can yield a uniform
lower bound of order \(t^{pq}\), let alone \(t^q\).
\end{proposition}

\begin{proof}
For \(\eps_n\le s\), one has
\(w_n\le s^{p-1}w_n\eps_n^{-(p-1)}\).  Hence
\[
 \frac{W(s)}{s^{p-1}}
 \le \sum_{\eps_n\le s}w_n\eps_n^{-(p-1)}.
\]
The right-hand side tends to zero as \(s\downarrow0\), because it is the
tail over those \(n\) for which \(\eps_n\) is small of the convergent series
in \eqref{eq:summability-barrier}.  This proves \eqref{eq:little-o}; setting
\(s=t^q\) gives \eqref{eq:barrier-gain}.
\end{proof}

Proposition~\ref{prop:barrier} is a limitation of this particular method, not
a counterexample to exact simultaneous renorming.  It explains why merely
optimizing the weights in the submitted construction cannot settle the
endpoint question.

\section{The exact endpoint for separable reflexive spaces}

The exact separable case is already contained, implicitly, in the tree/FDD
machinery of Odell and Schlumprecht \cite{OS}.  We spell out the deduction
because it is very close to the question in Remark~4.4 but does not appear to
have been stated there in this form.

\subsection{The order of the exponents}

\begin{lemma}\label{lem:p-le-q}
If an infinite-dimensional reflexive Banach space admits a \(p\)-AUS
renorming and a \(q\)-AUC renorming, then \(p\le q\).
\end{lemma}

\begin{proof}
Choose a normalized weakly null basic sequence.  A \(q\)-AUC norm yields,
after passage to a subsequence, a lower \(\ell_q\) estimate; a \(p\)-AUS
norm yields, after a further passage to a subsequence, an upper \(\ell_p\)
estimate.  These are the sequential special cases of
\cite[Proposition~3.3]{OS}.  Equivalent norms only change the constants.
Thus for some \(c,C>0\), the final subsequence \((x_i)\) satisfies
\[
 c\left(\sum|a_i|^q\right)^{1/q}
 \le \norm{\sum a_ix_i}
 \le C\left(\sum|a_i|^p\right)^{1/p}.
\]
Taking \(a_1=\cdots=a_n=1\) gives
\(cn^{1/q}\le Cn^{1/p}\) for every \(n\), hence \(p\le q\).
\end{proof}

\subsection{Tree estimates and exact FDD estimates}

Odell--Schlumprecht say that \(X\) satisfies \(C\)-\((P,Q)\)-tree estimates
if every normalized weakly null tree has one branch with a lower
\(\ell_P\) estimate and (possibly a different) branch with an upper
\(\ell_Q\) estimate.  The permission to use two different branches is the
key point.

Their Proposition~3.3 gives:
\begin{align*}
 q\text{-AUC}&\quad\Longrightarrow\quad(q,1)\text{-tree estimates},\\
 p\text{-AUS}&\quad\Longrightarrow\quad(\infty,p)\text{-tree estimates}.
\end{align*}
After fixing either original norm as a reference, normalize the nodes in the
other norm before applying Proposition~3.3 and then use equivalence of the two
norms to return to the reference normalization.  Thus the two separate
renormings give \((q,p)\)-tree estimates in one fixed reference norm.
Lemma~\ref{lem:p-le-q} supplies the required order
\(p\le q\).

The exact form of \cite[Theorem~4.1(e)]{OS} says that a separable reflexive
space with \((P,Q)\)-tree estimates embeds into a reflexive space having an
FDD with exact \(1\)-\((P,Q)\) estimates.  Applying it with
\((P,Q)=(q,p)\) gives a space \(Z\) with an FDD \((E_i)\) such that every
finite block sequence \((z_i)\) satisfies
\begin{equation}\label{eq:exact-fdd}
 \left(\sum_i\norm{z_i}^q\right)^{1/q}
 \le \norm{\sum_i z_i}
 \le \left(\sum_i\norm{z_i}^p\right)^{1/p}.
\end{equation}

\begin{lemma}\label{lem:fdd-moduli}
If a Banach space \(Z\) has an FDD satisfying \eqref{eq:exact-fdd}, then its
norm is \(p\)-AUS and \(q\)-AUC.
\end{lemma}

\begin{proof}
Fix \(x\in S_Z\) and \(0<t\le1\).  Given \(\eta>0\), choose a finitely
supported vector \(u\) with \(\norm{x-u}<\eta\), and let \(Y\) be the tail
of the FDD after the support of \(u\).  For \(y\in S_Y\), the two vectors
\(u\) and \(ty\) form a block sequence, so
\begin{align*}
 \norm{x+ty}
 &\le \eta+\bigl(\norm{u}^p+t^p\bigr)^{1/p},\\
 \norm{x+ty}
 &\ge \bigl(\norm{u}^q+t^q\bigr)^{1/q}-\eta.
\end{align*}
In the first inequality choose \(\eta\) to be a sufficiently small multiple
of \(t^p\); in the second choose it to be a sufficiently small multiple of
\(t^q\).  Since \(1-\eta\le\norm{u}\le1+\eta\), elementary estimates for
\((a^s+t^s)^{1/s}\) yield constants \(C,c>0\), independent of \(x,t\),
such that
\[
 \sup_{y\in S_Y}\norm{x+ty}\le1+Ct^p,
 \qquad
 \inf_{y\in S_Y}\norm{x+ty}\ge1+ct^q.
\]
This is precisely \(p\)-AUS and \(q\)-AUC.
\end{proof}

\begin{theorem}[Exact separable preservation]\label{thm:separable}
Let \(X\) be a separable reflexive Banach space.  If \(X\) admits separately
an equivalent \(p\)-AUS norm and an equivalent \(q\)-AUC norm, then \(X\)
admits one equivalent norm which is simultaneously \(p\)-AUS and
\(q\)-AUC.
\end{theorem}

\begin{proof}
The finite-dimensional case is immediate.  Assume that \(X\) is
infinite-dimensional.  As above, \(X\) satisfies \((q,p)\)-tree estimates.  By
\cite[Theorem~4.1(e)]{OS}, it embeds into a reflexive \(Z\) with an FDD
satisfying exact \(1\)-\((q,p)\) estimates.  Lemma~\ref{lem:fdd-moduli}
shows that the norm of \(Z\) is \(p\)-AUS and \(q\)-AUC.  Restricting this
norm to the embedded copy of \(X\) preserves both estimates.
\end{proof}

\begin{remark}
The paper \cite{OS} explicitly works with separable Banach spaces.  Thus
Theorem~\ref{thm:separable} does not by itself answer the nonseparable
question in \cite{DKLR}.  It does show that any obstruction to the general
endpoint must be genuinely nonseparable, rather than a phenomenon already
visible on one separable subspace.
\end{remark}

\section{Literature search and status of the nonseparable problem}

The literature search used the exact wording of Remark~4.4, variants of
``preserve both power types,'' citation searches around \cite{DKLR}, and
recent surveys and renorming papers.  The main findings are:

\begin{enumerate}[label=(\arabic*)]
\item The exact separable result follows from \cite[Proposition~3.3 and
Theorem~4.1(e)]{OS}, as proved in Theorem~\ref{thm:separable}.  The crucial
lines are the exact tree-to-FDD exponent transfer and the fact that lower and
upper branches need not coincide.
\item Lancien's 2017 lecture notes still list preservation of both exponents
as a question and emphasize that the fixed-power conditions are not suitable
for the elementary Baire-category proof used for qualitative simultaneous
AUS/AUC renorming \cite{LancienSlides}.
\item Recent work proves strong one-sided, separably determined renorming
theorems for power-type AUS properties; see, for example,
Causey--Fovelle--Lancien \cite{CFL}.  These results do not, as stated, produce
one global norm simultaneously realizing a prescribed AUS exponent and a
prescribed AUC exponent.
\item A 2024 paper of Baudier--Lancien describes the qualitative
simultaneous AUS/AUC renorming theorem as the current route to property
\((\beta)\), without providing an exact prescribed-power theorem
\cite{BLLocalGlobal}.
\item The 2025 monograph of Baudier--Lancien \cite{BLbook} provides a current
large-scale account of asymptotic renorming and the qualitative simultaneous
AUS/AUC problem.  I did not locate in it, or in papers citing \cite{DKLR}, a
published theorem or counterexample settling the exact endpoint for arbitrary
nonseparable reflexive spaces.
\end{enumerate}

The last statement is a report of the search, not a logical proof that no
resolution exists.  The safe current conclusion is therefore:

\begin{quote}
The exact endpoint is settled positively for separable reflexive spaces.  For
the arbitrary nonseparable setting of Remark~4.4, no published full answer or
counterexample was found in this search.  The positive linear Asplund
averaging method cannot reach the endpoint, and its correct general output is
Theorem~\ref{thm:corrected}.
\end{quote}

\section{Further directions}

The analysis isolates two plausible routes beyond the present result.

First, one could seek a nonseparable analogue of the exact tree-to-FDD
coordinatization theorem.  The obstacle is that asymptotic moduli in a
nonseparable space are naturally controlled by weakly null nets and
finite-codimensional subspaces, while the Odell--Schlumprecht construction is
sequential and explicitly separable.

Second, a successful averaging construction must be nonlinear in a substantial
way.  Proposition~\ref{prop:barrier} rules out obtaining the endpoint by merely
changing the positive weights or the sequence \((\eps_n)\) in a linear sum of
the predual norms of \(N^*+\eps_nM^*\).  Supremum-type, infimal-convolution,
or scale-adaptive constructions would have to avoid both the summability
condition \eqref{eq:summability-barrier} and the loss of AUC caused by a fixed
mixing parameter.

\section*{Acknowledgment of scope}

This note gives a complete audit of the submitted proof, a rigorous corrected
quantitative theorem for arbitrary reflexive spaces, a method barrier, and a
complete positive solution of the exact problem in the separable case.  It does
not claim to settle the remaining arbitrary nonseparable endpoint.

\begin{thebibliography}{99}

\bibitem{Asplund}
E. Asplund,
\newblock Averaged norms,
\newblock \emph{Israel J. Math.} 5 (1967), 227--233.

\bibitem{BLLocalGlobal}
F. P. Baudier and G. Lancien,
\newblock An asymptotic analog of a local-to-global phenomenon for uniformly convex renormings,
\newblock arXiv:2401.17465, 2024.

\bibitem{BLbook}
F. P. Baudier and G. Lancien,
\newblock \emph{Asymptotic and nonlinear geometries of Banach spaces and their interactions},
\newblock arXiv:2512.00817, 2025.

\bibitem{CFL}
R. M. Causey, A. Fovelle, and G. Lancien,
\newblock Asymptotic smoothness in Banach spaces, three space properties and applications,
\newblock \emph{Trans. Amer. Math. Soc.} 376 (2023), 1895--1928.

\bibitem{DGZ}
R. Deville, G. Godefroy, and V. Zizler,
\newblock \emph{Smoothness and Renormings in Banach Spaces},
\newblock Pitman Monographs and Surveys in Pure and Applied Mathematics 64,
Longman, 1993.

\bibitem{DKLR}
S. J. Dilworth, D. Kutzarova, G. Lancien, and N. L. Randrianarivony,
\newblock Equivalent norms with the property $(\beta)$ of Rolewicz,
\newblock \emph{Rev. R. Acad. Cienc. Exactas F\'is. Nat. Ser. A Mat. RACSAM}
111 (2017), 101--113; arXiv:1506.07978.

\bibitem{JLPS}
W. B. Johnson, J. Lindenstrauss, D. Preiss, and G. Schechtman,
\newblock Almost Fr\'echet differentiability of Lipschitz mappings between infinite-dimensional Banach spaces,
\newblock \emph{Proc. London Math. Soc.} 84 (2002), 711--746.

\bibitem{JohnZizler}
K. John and V. Zizler,
\newblock A short proof of a version of Asplund averaging theorem,
\newblock \emph{Proc. Amer. Math. Soc.} 73 (1979), 277--278.

\bibitem{LancienSlides}
G. Lancien,
\newblock \emph{A few questions on nonlinear embeddings into Banach spaces},
\newblock lecture notes, Svratka, January 2017.

\bibitem{OS}
E. Odell and Th. Schlumprecht,
\newblock Embedding into Banach spaces with finite dimensional decompositions,
\newblock \emph{Rev. R. Acad. Cienc. Exactas F\'is. Nat. Ser. A Mat. RACSAM}
100 (2006), 295--323; arXiv:math/0701324.

\end{thebibliography}

\end{document}
